\chapter{Fazit und Ausblick}
\label{ch:FazitundAusblick}

\todo[inline]{Fazit und Ausblick muss noch erstell werden.}

\section{Zusammenfassung}
\label{sec:Zusammenfassung}


\section{Ausblick}
\label{sec:ausblick}

Die vorliegende Arbeit hat gezeigt, dass eine robuste \ac{IBVS}-Regelung auch unter schwierigen Sichtbedingungen durch die Kombination aus Referenzmerkmalen, modellbasierter Vorhersage und Filterung realisierbar ist. Gleichzeitig ergeben sich aus den gewonnenen Erkenntnissen mehrere Ansätze für weiterführende Arbeiten, die das System sowohl hinsichtlich seiner Robustheit als auch hinsichtlich seiner praktischen Einsetzbarkeit weiter verbessern könnten.

Ein erster vielversprechender Ansatz besteht in der expliziten Berücksichtigung von Sichtfeldbeschränkungen innerhalb des Reglers. Im aktuellen System kann es in ungünstigen Bewegungsphasen dazu kommen, dass relevante Merkmale an den Rand des Bildes wandern oder das Sichtfeld vollständig verlassen. Zukünftige Arbeiten könnten daher eine Regelstrategie integrieren, die die Sichtbarkeit der verwendeten Features direkt als Nebenbedingung berücksichtigt. Denkbar wären hierbei beispielsweise Switching-Strategien, barrier-basierte Ansätze oder Verfahren mit vorgeschriebener Leistungsgrenze, welche die Bildtrajektorien gezielt innerhalb des Kamerasichtfelds halten. Dadurch könnte die Merkmalsverfügbarkeit während der gesamten Bewegung verbessert und die Stabilität des Regelkreises weiter erhöht werden. \cite{chesi2004keeping, heshmati2014ppibvs}

Ein zweiter relevanter Ausblick betrifft die automatische Kalibrierung der Filterparameter. In der vorliegenden Arbeit werden Prozess- und Messrauschen sowie weitere Schwellenwerte weitgehend manuell gewählt. Dies ist für die experimentelle Untersuchung sinnvoll, erschwert jedoch die Übertragbarkeit auf andere Objekte, Kamerakonfigurationen oder Bewegungsprofile. Eine zukünftige Erweiterung könnte daher adaptive Schätzverfahren einsetzen, die Kovarianzmatrizen und Gating-Parameter online aus den Innovationsgrößen oder Residuen des Filters ableiten. Auf diese Weise ließe sich die Parametrierung stärker an die aktuelle Beobachtungslage anpassen und die Gefahr einer Filterdivergenz bei wechselnden Sichtbedingungen reduzieren. \cite{birjandi2023adaptiveekf}

Ein weiterer Ansatz ist die explizite Kompensation von Systemlatenzen. Im praktischen Betrieb entstehen Verzögerungen durch Bildaufnahme, Keypoint-Erkennung, Matching, Filterung und die anschließende Berechnung der Stellgröße. Dadurch basiert das Regelsignal stets auf einem zeitlich verzögerten Bildzustand. Künftige Arbeiten könnten daher eine prädiktive beziehungsweise diskretzeitige \ac{IBVS}-Regelung integrieren, bei der nicht die zuletzt gemessene, sondern die für den tatsächlichen Wirkmoment vorhergesagte Merkmalslage zur Berechnung der Stellgröße verwendet wird. Da im entwickelten System bereits eine modellbasierte Keypoint-Vorhersage implementiert ist, bietet sich hier eine direkte Weiterentwicklung an. Insbesondere bei höheren Geschwindigkeiten oder schwankenden Rechenzeiten ist zu erwarten, dass eine solche Verzögerungskompensation die Regelgüte und Laufruhe weiter verbessert. \cite{yang2022sampleddata}

Darüber hinaus stellt die Erweiterung auf Mehrkamerasysteme beziehungsweise Sensor-Fallbacks einen interessanten Ausblick dar. Die aktuelle Arbeit basiert auf einer einzelnen Eye-in-Hand-Kamera, wodurch das System in besonderem Maß von deren Sichtfeld und Bildqualität abhängt. Zukünftig könnte eine zusätzliche Eye-to-Hand-Kamera oder ein hybrides Mehrkamerasystem eingebunden werden, um Sichtfeldverluste, Okklusionen oder ungünstige Perspektiven besser auszugleichen. Insbesondere eine Kombination aus lokaler Eye-in-Hand-Sicht für präzise Regelung und globaler Eye-to-Hand-Sicht für Sichtbarkeitsüberwachung und Relokalisierung erscheint hierfür vielversprechend. \cite{lippiello2005multicamera}

Schließlich bietet auch ein hybrider Visual-Servoing-Ansatz im Sinne des 2-1/2-D Visual Servoing ein interessantes Entwicklungspotenzial. Solche Verfahren kombinieren Eigenschaften des klassischen \ac{IBVS} mit zusätzlichen lagebezogenen Informationen, ohne die Nachteile eines vollständig posebasierten Reglers in gleichem Maße zu übernehmen. Für das in dieser Arbeit behandelte Szenario könnte ein hybrider Ansatz dazu beitragen, die Vorteile der direkten Bildregelung beizubehalten und gleichzeitig problematische Bewegungsphasen, schwache Beobachtbarkeit einzelner Freiheitsgrade oder ungünstige Trajektorien besser zu handhaben. Insbesondere die Verbindung aus robustem Keypoint-basiertem \ac{IBVS} und partieller geometrischer Information erscheint hierfür als sinnvoller nächster Schritt. \cite{malis1999twoandhalf}

Zusammenfassend zeigen diese Ansätze, dass die entwickelte Architektur eine geeignete Grundlage für weiterführende Forschungsarbeiten bildet. Insbesondere die Kombination aus sichtfeldbewusster Regelung, adaptiver Filterparametrierung, prädiktiver Latenzkompensation, Mehrkameraintegration und hybriden Visual-Servoing-Konzepten bietet ein hohes Potenzial, die Robustheit und industrielle Einsatzfähigkeit des Systems weiter zu steigern.

%% End of File %%